Sofern nicht explizit anders gefordert, sollten Arbeiten beidseitig gedruckt
werden, um Papier zu sparen. Hierzu ist die Klassenoption \texttt{twoside} zu
setzen. Für kürzere Haus- oder Seminararbeiten ist im Zweifel auch einseitiger
Druck statthaft.

\lstinputlisting[language={[LaTeX]{TeX}},style=block,escapechar=\!,caption={Doppelseitiger Textsatz},label={lst:tpl:twoside}]{code/tpl_twoside}

In umfangereicheren Arbeiten (ab Bachelorarbeit aufwärts), sollten neue
Kapitel stets auf einer rechten Seiten beginnen. Die Klasse \texttt{book} wird
vor Beginn eines neuen Kapitels im Bedarfsfall automatisch eine Leerseite
einfügen. Bei der Klasse \texttt{report} ist hierfür die Option
\texttt{openright} notwendig.

\lstinputlisting[language={[LaTeX]{TeX}},style=block,escapechar=\!,caption={Kapitelbeginn auf rechter Seite},label={lst:tpl:openright}]{code/tpl_openright}
